\chapter{Implementação e Stack Tecnológica}
\label{cap:implementacao}

Este capítulo trata da engenharia de \textit{software} da plataforma: como ela foi desenvolvida, quais tecnologias foram escolhidas e como as camadas cliente e servidor foram construídas. O objetivo é unir uma simulação local, rápida e fiel, a um servidor concorrente, seguro e de alto desempenho.

Enquanto o Capítulo~\ref{cap:arquitetura} foi organizado segundo o ciclo de vida da prova, este capítulo segue as camadas técnicas do sistema, pois cada componente atende a mais de uma fase. O Quadro~\ref{qua:fases-componentes} indica, para cada fase da prova, onde estão os componentes que a implementam.

\begin{quadro}[htbp]
    \caption{Correspondência entre as fases da prova e os componentes implementados}
    \label{qua:fases-componentes}
    \footnotesize
    \begin{tabularx}{\textwidth}{|>{\raggedright\arraybackslash}p{3.0cm}|L|>{\raggedright\arraybackslash}p{2.6cm}|}
        \hline
        \textbf{Fase da prova} & \textbf{Componentes} & \textbf{Seções} \\
        \hline
        Preparação da turma (Seção~\ref{sec:gestao-academica}) & Modelos GORM de usuários, turmas e matrículas; importação de CSV e envio assíncrono de e-mails & \ref{sec:gorm-persistencia} e \ref{sec:filas-emails-assincronos} \\
        \hline
        Autoria e montagem da prova (Seção~\ref{sec:autoria-conteudo}) & Serializador do VFS para os cenários; MinIO para imagens e \texttt{.tar.gz}; \textit{player} do YouTube para os trechos de vídeo; catálogo de componentes interativos; bloco de teste executado pelo mesmo interpretador do simulador; \textit{seed} do conteúdo atual & \ref{sec:vfs-typescript} e \ref{sec:gorm-persistencia} \\
        \hline
        Abertura da prova (Seção~\ref{sec:abertura-prova}) & Rotas de \textit{token} e janelas no Gin; autenticação JWT; sessão única no Hub WebSocket & \ref{sec:gin-gonic} e \ref{sec:websocket-hub} \\
        \hline
        Resolução (Seção~\ref{sec:validacao-desacoplada}) & Simulador POSIX e VFS no navegador; motor de asserções em \textit{worker pool} no servidor & \ref{sec:vfs-typescript} e \ref{sec:websocket-hub} \\
        \hline
        Supervisão (Seção~\ref{sec:governanca-seguranca}) & Janelas de ajuda e avisos no cliente; Hub WebSocket com eventos para o \textit{Cockpit} & \ref{sec:interface-ts} e \ref{sec:websocket-hub} \\
        \hline
        Encerramento (Seção~\ref{sec:encerramento-prova}) & Registro no IndexedDB e pacote \texttt{.proof}; transações e histórico preservado no PostgreSQL & \ref{sec:visao-stack} e \ref{sec:gorm-persistencia} \\
        \hline
    \end{tabularx}
    \fonte{Autoria própria (2026).}
\end{quadro}

% ----------------------------------------------------------
\section{Metodologia de Engenharia: Adoção Instrumental de Inteligência Artificial e Governança Humana}
\label{sec:metodologia-ia}
% ----------------------------------------------------------

A engenharia de \textit{software} passa por uma mudança estrutural com a chegada de assistentes de IA generativa e de Modelos de Linguagem de Grande Escala (LLMs) aos fluxos de desenvolvimento. Em empresas de tecnologia, o uso desses assistentes como ferramenta já se consolidou para acelerar a produção, prototipar com agilidade e refinar documentação técnica, sem que isso signifique transferir responsabilidade à máquina.

Com base em mais de cinco anos de experiência do autor em liderança técnica no mercado, este trabalho adota essa forma de desenvolvimento acelerado. O uso de IA segue o princípio do \textit{Human-in-the-Loop} (humano no circuito de controle): a ferramenta amplia a capacidade de execução, mas tudo passa por controle humano.

\begin{enumerate}
    \item \textbf{concepção e decisões de domínio:} a ideia do projeto, o modelo de dados, os protocolos criptográficos da contingência \textit{offline}, as regras disciplinares e a dinâmica das avaliações vieram do diagnóstico dos laboratórios da UTFPR e da experiência do autor e da orientadora. Nenhuma regra de negócio, limite de segurança ou decisão de arquitetura foi delegada à máquina;
    \item \textbf{aceleração e apoio à escrita:} o assistente foi usado para automatizar tarefas repetitivas, sugerir rotinas de infraestrutura, transformar conceitos em especificações e ajustar a redação técnica às normas acadêmicas;
    \item \textbf{revisão crítica:} como discutido na Seção~\ref{sec:risco-ia}, o valor do profissional de tecnologia hoje está em auditar o que a ferramenta produz. Cada parágrafo, comando simulado e contrato de API desta monografia foi revisado, testado e validado pelo autor.
\end{enumerate}

Assim, embora a forma do texto e o ritmo de trabalho tenham se beneficiado dessas ferramentas, a concepção, a arquitetura e a responsabilidade científica do trabalho são inteiramente humanas, como já acontece nas práticas mais avançadas da indústria.

% ----------------------------------------------------------
\section{Visão Geral da Stack Tecnológica}
\label{sec:visao-stack}
% ----------------------------------------------------------

As tecnologias foram escolhidas por quatro critérios: escalabilidade, fidelidade da simulação, facilidade de implantação na infraestrutura da universidade e separação clara de responsabilidades. O Quadro~\ref{qua:stack} resume as escolhas.

\begin{quadro}[htbp]
    \caption{Tecnologias adotadas em cada camada}
    \label{qua:stack}
    \footnotesize
    \begin{tabularx}{\textwidth}{|p{2.2cm}|p{3.4cm}|L|}
        \hline
        \textbf{Camada} & \textbf{Tecnologia} & \textbf{Por que foi escolhida} \\
        \hline
        Cliente & TypeScript & Tipagem estática para modelar com precisão as estruturas do POSIX (nós, permissões, credenciais) \\
        \hline
        Cliente & DOM nativo, sem \textit{framework} & Telas escritas como classes TypeScript que montam e desmontam seus próprios elementos, sem dependências externas de interface \\
        \hline
        Cliente & Web Crypto API e IndexedDB & \textit{Hashes} e assinaturas nativos do navegador e armazenamento local transacional de todo o registro da prova, base da contingência \textit{offline} \\
        \hline
        Servidor & Go & Concorrência leve com \textit{goroutines} e canais, ideal para centenas de conexões simultâneas \\
        \hline
        Servidor & Gin-Gonic & Roteamento HTTP rápido, com \textit{middlewares} para autenticação e autorização \\
        \hline
        Servidor & GORM e PostgreSQL & Mapeamento objeto-relacional tipado, transações ACID e exclusão lógica \\
        \hline
        Cliente & YouTube IFrame Player API & Exibição dos trechos de vídeo dos blocos, com início e fim, e leitura da duração para validar o trecho no editor \\
        \hline
        Servidor & MinIO & Armazenamento de objetos para as imagens dos módulos e os arquivos \texttt{.tar.gz} dos cenários \\
        \hline
        Comunicação & WebSockets (WSS) & Canal bidirecional persistente para telemetria e eventos em tempo real \\
        \hline
    \end{tabularx}
    \fonte{Autoria própria (2026).}
\end{quadro}

% ----------------------------------------------------------
\section{Front-end e Núcleo de Simulação (Client-side)}
\label{sec:frontend-simulacao}
% ----------------------------------------------------------

Rodar a simulação no navegador evita tanto as máquinas virtuais quanto um contêiner Docker por aluno, o que seria inviável para laboratórios com orçamento de infraestrutura limitado.

\subsection{Modelagem do VFS e Execução POSIX em TypeScript}
\label{sec:vfs-typescript}

O núcleo da simulação foi escrito inteiramente em TypeScript. A tipagem estática é essencial para representar com precisão as estruturas do POSIX:

\begin{enumerate}
    \item \textbf{nós e metadados (\textit{i-nodes}):} cada entrada do VFS é modelada por interfaces que representam arquivos regulares, diretórios e \textit{links} simbólicos, com permissões em notação octal e simbólica, dono e grupo (UID/GID), referências aos filhos e datas de modificação;
    \item \textbf{permissões e \texttt{umask}:} as permissões \textit{rwx} e a máscara de criação \texttt{umask} são calculadas com operações bit a bit (\textit{bitwise}), reproduzindo o controle de acesso discricionário (DAC) do Linux;
    \item \textbf{troca de usuário:} o simulador mantém uma pilha de sessões que reproduz a troca de credenciais com \texttt{su} e \texttt{sudo}, incluindo as restrições do arquivo \texttt{/etc/sudoers};
    \item \textbf{interpretador de comandos:} o \textit{parser} trata \textit{pipelines} (\texttt{|}), redirecionamentos (\texttt{>}, \texttt{>>}, \texttt{<}) e expansão de variáveis de ambiente, tudo no navegador e sem atraso na digitação.
\end{enumerate}

\subsection{Componentização da Interface em TypeScript}
\label{sec:interface-ts}

A interface é escrita em TypeScript puro, manipulando diretamente o DOM (\textit{Document Object Model}) do navegador, sem \textit{framework} de interface. Cada tela é uma classe que sabe se montar em um contêiner e se desmontar ao sair, e a navegação entre telas é feita pelo endereço da página. Os principais módulos são:

\begin{enumerate}
    \item \textbf{terminal ANSI:} emulador compatível com sequências de escape ANSI e com a paleta do GNOME Terminal, com autocompletar por \texttt{Tab} para comandos e caminhos, histórico pelas setas e os editores \texttt{nano} e \texttt{vim};
    \item \textbf{painel de roteiros e telemetria discente:} painel retrátil com objetivos, enunciados, progresso de questões e cronômetro sincronizado periodicamente com o relógio do servidor;
    \item \textbf{entrada da prova e contingência:} tela para validação do \textit{token} emitido pela professora e captura de pacotes criptográficos em modo desconectado;
    \item \textbf{grade de supervisão do Cockpit docente:} componente reativo em mosaico que renderiza os \textit{cards} discentes atualizados via \texttt{cockpit:telemetry-update}, ordenando dinamicamente a fila de atendimento prioritário (``mãozinha virtual'') e exibindo métricas de tempo de permanência por questão;
    \item \textbf{espelhamento ao vivo (\textit{Live Terminal Mirror}):} painel modal do \textit{Cockpit} que renderiza uma réplica visual síncrona do terminal do aluno e da árvore de arquivos do VFS, alimentado pelo evento \texttt{cockpit:terminal-mirror};
    \item \textbf{janelas de apoio e comunicação direta:} modais para solicitação de ajuda com justificativa obrigatória, telas de advertência escalonada na detecção de \textit{DevTools} e camada sobreposta translúcida para recepção de orientações privadas da professora (\texttt{teacher:direct-message}).
\end{enumerate}

\subsection{Funcionalidades Utilitárias de Perfil}
\label{sec:perfil-utilitario}

A interface também inclui as telas de autoedição de perfil, comuns a alunos e professores: dados de contato, avatar e redefinição de senha com o usuário autenticado. Trata-se de formulários e rotas CRUD convencionais, sem particularidades em relação à prática corrente da indústria, e por isso não são aprofundados aqui. Como explicado na Seção~\ref{sec:modelo-dados}, o detalhamento técnico deste capítulo prioriza o que diferencia o projeto: o núcleo de emulação POSIX e o VFS, a mitigação de fraudes via \textit{DevTools}, a telemetria via \textit{WebSockets} e a resiliência a falhas de rede.

% ----------------------------------------------------------
\section{Back-end e Concorrência de Alta Performance (Server-side)}
\label{sec:backend-golang}
% ----------------------------------------------------------

O servidor é a autoridade confiável da plataforma: guarda as regras de negócio, os dados acadêmicos e garante a integridade das avaliações.

\subsection{A Linguagem Golang e o Modelo de Concorrência CSP}
\label{sec:golang-concorrencia}

A linguagem Go \cite{donovan2015go} foi escolhida pela capacidade de lidar com muita concorrência de rede gastando pouca memória \cite{cox2020golang}:

\begin{enumerate}
    \item \textbf{\textit{goroutines} leves:} \textit{threads} do sistema operacional reservam de 1 a 2\,MB de pilha cada uma, e modelos de \textit{loop} de eventos em uma única \textit{thread} podem travar durante a correção de asserções. Já as \textit{goroutines} começam com pilhas de apenas 2\,KB, que crescem conforme a necessidade. Isso permite manter centenas de conexões \textit{WebSocket} por sala com impacto mínimo na memória do servidor;
    \item \textbf{canais (CSP):} telemetria, fila de dúvidas e despacho de eventos trocam mensagens por canais tipados (\textit{channels}), inspirados no modelo de Processos Sequenciais Comunicantes (\textit{Communicating Sequential Processes} -- CSP) \cite{donovan2015go,cox2020golang}. Isso evita condições de corrida (\textit{race conditions}) e reduz a necessidade de travas manuais (\textit{mutexes}).
\end{enumerate}

\subsection{Roteamento de APIs REST com Gin-Gonic}
\label{sec:gin-gonic}

As APIs seguem o estilo arquitetural REST \cite{fielding2000rest} e são servidas pelo \textit{framework} \textbf{Gin-Gonic}, que usa árvores de prefixo (\textit{radix tree}) para rotear requisições com rapidez e sem alocações desnecessárias de memória:

\begin{enumerate}
    \item \textbf{autenticação e autorização:} \textit{middlewares} verificam os \textit{tokens} JWT, os papéis de usuário (RBAC: Público, Aluno e Professor) e a expiração das sessões;
    \item \textbf{rotas de negócio:} gestão de turmas e matrículas, criação de roteiros de avaliação, emissão e validação dos \textit{tokens} de sala e homologação de notas com justificativa do professor.
\end{enumerate}

\subsection{Persistência e Mapeamento Objeto-Relacional com GORM}
\label{sec:gorm-persistencia}

A persistência de dados utiliza a biblioteca \textbf{GORM} integrada ao PostgreSQL, combinando integridade referencial estrita e modelagem orientada a objetos em Go:

\begin{enumerate}
    \item \textbf{modelos tipados e migrações:} usuários, turmas, matrículas, avaliações, sessões, submissões, observações do professor e eventos da linha do tempo são \textit{structs} Go com \textit{tags} GORM que configuram chaves primárias UUID, índices compostos e integridade referencial;

    \item \textbf{extensão relacional de submissões e sessões:} a entidade de respostas individuais (\texttt{respostas\_questoes\_sessoes}) foi estendida com suporte completo ao fluxo de revisão formativa e auditoria forense. Cada registro possui identificador próprio, chaves estrangeiras para a sessão e a questão (\texttt{sessao\_id} e \texttt{questao\_id}), a pontuação preliminar calculada pelo motor avaliador (\texttt{nota\_sugerida}), a pontuação final confirmada ou revisada pela professora (\texttt{nota\_homologada}), a fundamentação textual para eventuais divergências (\texttt{justificativa\_ajuste}), a devolutiva pedagógica específica (\texttt{feedback\_docente}) e a trilha completa de comandos digitados, tempos de resposta e saídas em coluna JSONB (\texttt{comandos\_executados}, mapeada via \texttt{datatypes.JSON}). Complementarmente, a tabela de sessões de avaliação (\texttt{sessoes\_prova}) incorpora o campo global \texttt{parecer\_geral\_docente}, a nota consolidada (\texttt{nota\_final\_homologada}), o carimbo temporal da homologação (\texttt{homologado\_em}) e o identificador do docente responsável (\texttt{homologado\_por\_id});

    \item \textbf{homologação transacional ACID:} a validação, recálculo e homologação das notas operam estritamente sob transações ACID gerenciadas pelo GORM (\texttt{tx := db.Begin()} \ldots{} \texttt{tx.Commit()}). Dentro do escopo transacional, o serviço atualiza em lote as notas homologadas e os \textit{feedbacks} individuais por questão, recalcula atomicamente a média ou soma ponderada final da avaliação persistindo-a em \texttt{sessoes\_prova}, grava o parecer pedagógico geral e registra um evento irrevogável na linha do tempo de auditoria com o carimbo de tempo da ação. Qualquer inconsistência ou falha de rede aciona reversão integral (\texttt{tx.Rollback()}), impedindo visualizações parciais ou notas descompassadas no portal do estudante;

    \item \textbf{transferência de turma em uma transação:} a troca de turma roda igualmente em uma transação ACID: a matrícula de origem passa a \texttt{TRANSFERIDO} e a nova matrícula é criada ativa na turma de destino em uma única operação. Submissões, notas e comandos antigos continuam ligados ao UUID e ao RA do aluno, sem registros órfãos nem cadastros duplicados;

    \item \textbf{exclusão lógica com preservação de auditoria:} a entidade de usuário usa o tipo \texttt{gorm.DeletedAt}. Nas consultas comuns, o GORM acrescenta automaticamente a condição \texttt{deleted\_at IS NULL}. Para auditorias junto à coordenação da UTFPR, consultas com \texttt{db.Unscoped()} recuperam todo o histórico, inclusive os registros de alunos desligados;

    \item \textbf{suspensão com efeito imediato:} quando o professor marca um aluno como \texttt{INATIVO}, o serviço grava a mudança no banco e envia um sinal, por canal interno, ao Hub de \textit{WebSockets}. As conexões abertas daquele usuário são encerradas na hora, e as próximas requisições à API são barradas pelo \textit{middleware} de autenticação.
\end{enumerate}

\subsection{Hub WebSocket Nativo e Streaming de Telemetria}
\label{sec:websocket-hub}

A comunicação bidirecional em tempo real é sustentada por um \textbf{Hub de WebSockets} nativo desenvolvido em Go, em estrita conformidade com a RFC 6455 \cite{fette2011websocket}:

\begin{enumerate}
    \item \textbf{Hub central e roteamento por sala:} uma \textit{goroutine} central opera como despachante de eventos, gerenciando o registro de novos terminais (\texttt{register}), o encerramento seguro (\texttt{unregister}) e a distribuição de mensagens segmentadas por sala de avaliação;

    \item \textbf{arquitetura de duas \textit{goroutines} por conexão:} cada cliente conectado (discente ou docente) é atendido por um par independente de rotinas assíncronas: uma rotina de leitura contínua (\textit{readPump}) e uma de despacho com amortecimento (\textit{writePump}). Esse desacoplamento previne que eventuais latências de rede ou pausas de renderização no navegador de um estudante comprometam a entrega de pacotes aos demais pares da sala;

    \item \textbf{eventos dedicados de supervisão do Cockpit:} o protocolo define mensagens tipadas em formato JSON para operacionalizar a telemetria em tempo real:
    \begin{itemize}
        \item \texttt{cockpit:telemetry-update}: transporta cargas leves com o total de questões concluídas e pendentes, a identificação da questão corrente e as métricas de temporização (tempo decorrido na tarefa atual e média acumulada por questão);
        \item \texttt{cockpit:terminal-mirror}: transmite em fluxo contínuo os blocos incrementais de caracteres digitados no terminal, códigos de retorno, saídas textuais e o estado serializado da árvore do VFS;
        \item \texttt{teacher:direct-message}: despacha mensagens de intervenção pedagógica direta e orientações privativas (1:1) do docente para a conexão exclusiva do estudante;
    \end{itemize}

    \item \textbf{desacoplamento de goroutines e prevenção de gargalos de I/O no espelhamento:} o espelhamento contínuo da digitação de múltiplos alunos em tempo real possui potencial para saturar o barramento de rede caso implementado via difusão indiscriminada. Para eliminar qualquer risco de contenção de I/O no servidor:
    \begin{itemize}
        \item o Hub adota subscrição seletiva sob demanda (\textit{targeted subscription}): eventos \texttt{cockpit:terminal-mirror} só são despachados aos canais docentes que tenham solicitado ativamente a abertura da tela espelhada daquele discente específico no painel;
        \item as transmissões utilizam canais com \textit{buffer} calibrados para o tráfego em rajadas típico de digitação (\texttt{chan []byte});
        \item emprega-se seleção não bloqueante (\texttt{select} com cláusula \texttt{default}) nas filas de espelhamento transitório. Se o canal de saída de um observador atingir a capacidade máxima por congestionamento temporário de rede, quadros intermediários do espelho são graciosamente descartados sem bloquear a \textit{goroutine} de leitura do aluno nem degradar os canais críticos de notas e batimentos de integridade (\textit{heartbeats});
    \end{itemize}

    \item \textbf{avaliação desacoplada fora do canal de mensageria:} ao término de cada comando prático acompanhado de \textit{snapshot} do sistema de arquivos virtual, a reconstrução e a verificação de asserções são descarregadas para um conjunto de trabalhadores assíncronos (\textit{worker pool}). O canal WebSocket permanece desimpedido para novas interações, assegurando resposta com latência imperceptível para o aluno.
\end{enumerate}

\subsection{Processamento Assíncrono de Filas, Agendamento e Despacho Multicanal}
\label{sec:filas-emails-assincronos}

Importar uma turma por CSV, com dezenas ou centenas de alunos, ou publicar um aviso para a turma gera muitos e-mails. Enviá-los via SMTP dentro da própria requisição HTTP seria inviável: negociar a conexão TLS e enviar cada mensagem leva de 500 milissegundos a 2 segundos por destinatário, o que estouraria o tempo limite (\textit{Gateway Timeout}, HTTP 504) e deixaria a professora esperando.

Por isso, o servidor usa um modelo produtor-consumidor (\textit{producer-consumer}) com um conjunto fixo de trabalhadores (\textit{worker pool}) e um agendador, ilustrado na Figura~\ref{fig:worker-pool}.

\begin{figure}[htbp]
    \centering
    \caption{Envio assíncrono de e-mails com canal e conjunto de trabalhadores}
    \label{fig:worker-pool}
    \begin{tikzpicture}[node distance=0.9cm and 1.1cm]
        \node[box, text width=2.6cm] (csv) {Importação CSV\\(\textit{handler} Gin)};
        \node[box, text width=2.6cm, below=0.5cm of csv] (sched) {Agendador\\(\texttt{time.Ticker})};
        \node[box, text width=2.5cm, right=1.4cm of $(csv)!0.5!(sched)$] (chan) {Canal\\\texttt{chan EmailJob}};
        \node[box, text width=1.9cm, right=1.4cm of chan, yshift=1.0cm] (w1) {\textit{worker} 1};
        \node[box, text width=1.9cm, right=1.4cm of chan] (w2) {\textit{worker} 2};
        \node[box, text width=1.9cm, right=1.4cm of chan, yshift=-1.0cm] (wn) {\textit{worker} $N$};
        \node[box, text width=1.6cm, right=1.1cm of w2] (smtp) {Servidor\\SMTP};

        \draw[arrow] (csv.east) -- (chan);
        \draw[arrow] (sched.east) -- (chan);
        \draw[arrow] (chan) -- (w1.west);
        \draw[arrow] (chan) -- (w2.west);
        \draw[arrow] (chan) -- (wn.west);
        \draw[arrow] (w1.east) -- (smtp);
        \draw[arrow] (w2.east) -- (smtp);
        \draw[arrow] (wn.east) -- (smtp);
        \node[font=\scriptsize, below=0.3cm of sched] {responde \texttt{202 Accepted} na hora};
        \node[font=\scriptsize, below=0.3cm of wn] {limite de taxa e \textit{backoff}};
    \end{tikzpicture}
    \fonte{Autoria própria (2026).}
\end{figure}

O fluxo funciona assim:

\begin{enumerate}
    \item \textbf{fila em canal com \textit{buffer}:} o \textit{handler} HTTP valida os dados (CSV ou novo aviso), grava tudo no banco em uma transação e coloca as tarefas de envio em um canal tipado com \textit{buffer} (\texttt{chan EmailJob}). A resposta \texttt{202 Accepted} volta ao cliente em poucos milissegundos;
    \item \textbf{agendador de avisos:} para avisos programados, uma \textit{goroutine} com temporizador (\texttt{time.Ticker}) consulta periodicamente o PostgreSQL. Quando chega o horário de \texttt{publicar\_em}, o aviso passa a \texttt{PUBLICADO}, o evento \texttt{announcement:new} é enviado aos alunos conectados e os e-mails entram na fila;
    \item \textbf{trabalhadores com limite de taxa:} um número fixo $N$ de \textit{goroutines} consome o canal continuamente. Um limitador de taxa (\textit{rate limiter}) controla o ritmo de envio para respeitar as cotas e as regras antispam do servidor de e-mail da universidade;
    \item \textbf{e-mails em HTML:} cada trabalhador monta a mensagem a partir de modelos HTML com estilos \textit{inline}, compatíveis com os principais leitores de e-mail, com o logotipo da instituição, a identificação da disciplina, o texto do aviso e \textit{links} de ação protegidos por \textit{tokens} de uso único;
    \item \textbf{confirmação de leitura em tempo real:} quando o aluno confirma a leitura de um aviso importante, o servidor grava data, hora e IP e envia \texttt{cockpit:announcement-receipt}, atualizando na hora o percentual de adesão e a lista de pendentes no painel do professor;
    \item \textbf{novas tentativas com espera crescente:} se o servidor SMTP falhar temporariamente (sobrecarga ou erro no TLS), o trabalhador tenta de novo com intervalos cada vez maiores (\textit{exponential backoff}). Esgotadas as tentativas, o erro é registrado em \textit{log} e a notificação fica marcada como pendente, para que a professora possa reenviá-la pelo painel.
\end{enumerate}
